% Generated by tools/paper/render.py from main_text.py; edit the source there, not this file.
\section{Discussion}

sporos computes the seeds used by read mappers, genome aligners and pangenome tools, following the established definitions, and on a single core it was faster than each tool it can replace. Its frequency-capped MEMs show that how often each match's string occurs can be derived from the set of matches itself: counting containing intervals replaces an index of the genomes and reduces capped matching between two chromosomes to about a second.

A single library also avoids converting between tools. A pipeline loads a reference once and requests MEMs, MUMs or SMEMs from Mojo, C or Python, with one coordinate system and one set of occurrence limits, and the command-line tool writes the MUMmer and \texttt{bwa fastmap} formats that existing scripts expect.

\textbf{Limitations.} All measurements were taken on one x86-64 processor; the 8-lane rows of Table~\ref{tab:seeds} indicate performance at AVX2 width, but other processors, including ARM, were not benchmarked. Mumemto was compared on pairs of genomes rather than the large collections it targets, and its definitions differ from those of sporos. Positions are 32-bit, limiting each sequence and the concatenated reference to fewer than \(2^{32}\) bases, and FM-index construction needs about ten bytes of memory per reference base. On ten threads sporos was slower than simd-minimizers for one of the eight schemes.
